\chapter{Arquitectura del Sistema}
\label{ch:architecture}
% ==============================================================

\section{Stack tecnológico}

\begin{table}[H]
\centering
\caption{Dependencias principales}
\label{tab:stack}
\begin{tabular}{L{4cm} L{3.5cm} L{5cm}}
\toprule
\textbf{Componente} & \textbf{Versión} & \textbf{Rol} \\
\midrule
Python        & 3.11         & Runtime principal \\
FastAPI       & $\geq$ 0.110 & API REST backend \\
OpenCV        & 4.x          & Procesamiento de imagen \\
NumPy/SciPy   & --           & Álgebra matricial \\
SAM (Meta AI) & --           & Segmentación semántica \\
PyTorch       & $\geq$ 2.0   & Backend SAM \\
Vue 3 + Vite  & --           & Interfaz web \\
Tailwind CSS  & 3.x          & Estilos \\
fpdf2         & --           & Generación de informes PDF de calibración \\
PyInstaller + Inno Setup & -- & Empaquetado del instalador de escritorio \\
\bottomrule
\end{tabular}
\end{table}

\begin{infobox}[¿Por qué no hay GDAL ni pyproj en la lista?]
El sistema no reproyecta entre sistemas de referencia distintos en tiempo de ejecución: las
coordenadas UTM de las varillas de calibración ya vienen dadas en EPSG:25830 (capítulo~\ref{ch:calibration}),
así que la homografía píxel→UTM que calcula \texttt{calculate\_homography()} proyecta
directamente a ese sistema — \texttt{"EPSG:25830"} se escribe como metadato fijo en el GeoJSON
de salida (\texttt{georef\_export.py}), no se calcula con una librería de transformación de CRS.
Si en el futuro hiciera falta soportar más de un huso UTM o reproyectar a otro sistema, ahí sí
haría falta pyproj.
\end{infobox}

\newpage
\section{Estructura de directorios}

\begin{codebox}[Árbol del repositorio]
\begin{minted}{text}
cv-lit/
+-- acces_api/              # Módulo Obscape API
|   +-- obscape_api.py      # Cliente REST
|   +-- scheduled_download.py
+-- backend/                # Servidor FastAPI
|   +-- main.py             # App principal + endpoints de cámaras/calibración
|   +-- batch_alignment.py  # Pipeline de alineación masiva
|   +-- auto_mode.py        # Auto Mode: descarga + alineación + análisis
|   +-- config.py           # Constantes globales + CAMERAS
|   +-- desktop_launcher.py # Punto de entrada de la app de escritorio (pywebview)
|   +-- desktop_launcher.spec # Spec de PyInstaller para el .exe principal
|   +-- download_sam.py     # Descargador aparte del checkpoint SAM (opcional)
|   +-- download_sam.spec   # Spec de PyInstaller para download_sam.exe
+-- installer/               # Instalador de escritorio para Windows
|   +-- cv-lit.iss          # Script de Inno Setup
|   +-- output/             # LineaDeCosta-Setup.exe (generado, excl. Git)
+-- calibration/            # Perfiles y máscaras por cámara
|   +-- cam_N_H.npy         # Matriz H 3×3 en float64
|   +-- cam_N_profile.json  # Metadatos, RMSE, nombre de perfil, operador, notas
|   +-- roi_config.json     # ROIs por cámara
|   +-- alignment_masks.json# Zonas SIFT estables
+-- data/                   # Imágenes y resultados (excl. Git) — CVLIT_DATA_DIR
|   +-- camera1/ ... camera6/  # una carpeta por cámara (ver CAMERAS[i]["folder"])
|   |   +-- *.jpg           # capturas originales, nombradas <epoch>_<AAAAMMDD>_<HHMMSS>_<serie>.jpg
|   |   +-- aligned/        # salida del commit de alineación masiva (batch_alignment.py)
|   +-- CAM_1/ ... CAM_6/
|   |   +-- json/           # anotaciones de varillas (una por imagen marcada)
|   +-- system_logs.jsonl   # buffer de actividad que sirve GET /api/logs
+-- docs/                   # Documentación técnica
|   +-- memoria.tex         # Esta memoria  [NUEVO]
|   +-- reports/
+-- frontend/               # Interfaz Vue 3
|   +-- src/
|       +-- App.vue         # Layout raíz + navegación
|       +-- api.js          # API_BASE (VITE_API_URL) usado por todos los componentes
|       +-- components/
|           +-- Dashboard.vue        # Vista general del sistema
|           +-- ImageIngest.vue      # Carga de imágenes
|           +-- Calibration.vue      # Wizard de calibración (7 pasos)
|           +-- BatchAlignment.vue   # Alineación masiva
|           +-- AutoMode.vue         # Auto Mode
|           +-- CoastlineAnalysis.vue# Resultados de análisis
|           +-- GeoJSONMap.vue       # Mapa GeoJSON exportado
|           +-- Map.vue              # Mapa base reutilizable
|           +-- Cameras.vue          # Gestión de cámaras
+-- homography/
|   +-- align_images.py     # CLI standalone de alineación
+-- proces_images/          # Algoritmos de procesamiento
|   +-- extract_coastline.py
|   +-- segmentation_sam.py
|   +-- georef_export.py
+-- scripts/                # Herramientas de operación
+-- tests/
\end{minted}
\end{codebox}

\section{Pipeline de procesamiento}

La figura~\ref{fig:pipeline} muestra el grafo completo de procesamiento desde la adquisición
hasta la exportación del GeoJSON georreferenciado.

\begin{figure}[H]
\centering
\begin{tikzpicture}[node distance=0.8cm and 1.5cm]

  % Nodos
  \node[block, fill=Teal!12, draw=Teal!70] (api)
    {API Obscape\\{\tiny Descarga imágenes (14d)}};

  \node[block, fill=Orange!10, draw=Orange!70, below=of api] (batch)
    {\textbf{[3.0b] Alineación masiva}\\{\tiny SIFT + RANSAC · commit no-destructivo}};

  \node[block, below=of batch] (preproc)
    {[3.1] Preprocesado\\{\tiny Normalización + recorte ROI}};

  \node[block, below=of preproc] (sam)
    {[3.2] Segmentación SAM\\{\tiny Máscara arena/agua}};

  \node[block, below=of sam] (extract)
    {[3.3] Extracción línea de costa\\{\tiny cv2.findContours}};

  \node[block, below=of extract] (proj)
    {[3.4] Proyección homografía\\{\tiny $(u,v) \rightarrow$ EPSG:25830}};

  \node[block, fill=OliveGreen!10, draw=OliveGreen!70, below=of proj] (export)
    {[3.5] Exportación GeoJSON\\{\tiny EPSG:25830 · área seca m²}};

  \node[block, fill=NavyBlue!8, draw=NavyBlue!50, right=2.5cm of batch] (calib)
    {[3.0] Calibración offline\\{\tiny GCPs GNSS · RANSAC · H por cámara}};

  % Flechas verticales
  \draw[arrow] (api) -- (batch);
  \draw[arrow] (batch) -- (preproc);
  \draw[arrow] (preproc) -- (sam);
  \draw[arrow] (sam) -- (extract);
  \draw[arrow] (extract) -- (proj);
  \draw[arrow] (proj) -- (export);

  % Calibración lateral
  \draw[arrow, dashed, color=NavyBlue!50] (calib) -- (proj)
    node[midway, right, label] {H\textsubscript{i}};
  \draw[arrow, dashed, color=NavyBlue!50] (calib) -- (batch)
    node[midway, right, label] {máscaras};

\end{tikzpicture}
\caption{Pipeline completo de procesamiento cv-lit}
\label{fig:pipeline}
\end{figure}

% ==============================================================
